\section{Khảo sát sơ bộ trên 15 loại MVTec AD}
\label{sec:khao-sat}

Trước khi chọn các loại sản phẩm cho thực nghiệm chính, nhóm chạy cả ba phương pháp trên đủ 15 loại của MVTec AD
theo cùng một giao thức: ảnh $256\times256$, chia train/test gốc, seed 42, một lần chạy trên GPU Tesla T4. Khảo sát có
hai mục đích: kiểm tra cách đánh giá của nhóm có tái lập được kết quả đã công bố hay không, và tìm các loại làm lộ
điểm yếu của từng phương pháp. Ở bước này Autoencoder là mạng của nhóm, còn PaDiM và PatchCore được chạy bằng thư viện
Anomalib~\cite{akcayAnomalibDeepLearning2022}, với PaDiM dùng backbone ResNet-18 và PatchCore dùng Wide-ResNet-50-2.
Mọi kết quả từ Mục~\ref{sec:ket-qua-tong-hop} trở đi dùng mã tự cài trong \texttt{src/tn} và cùng backbone
Wide-ResNet-50-2 cho cả PaDiM và PatchCore.

\begin{table}[htbp]
  \centering
  \caption{AUROC trung bình trên 15 loại MVTec AD trong khảo sát sơ bộ}
  \label{tab:khao-sat}
  \begin{tabular}{@{}llccc@{}}
    \toprule
    Chỉ số & Nhóm loại & Autoencoder & PaDiM & PatchCore \\
    \midrule
    Image-AUROC & Bề mặt (5 loại)  & 0,784 & 0,945 & 0,993 \\
                & Vật thể (10 loại) & 0,741 & 0,863 & 0,977 \\
                & Cả 15 loại       & 0,756 & 0,890 & 0,983 \\
    \midrule
    Pixel-AUROC & Bề mặt (5 loại)  & 0,781 & 0,944 & 0,970 \\
                & Vật thể (10 loại) & 0,875 & 0,968 & 0,985 \\
                & Cả 15 loại       & 0,843 & 0,960 & 0,980 \\
    \bottomrule
  \end{tabular}
\end{table}

PatchCore cao nhất ở 13/15 loại theo image-AUROC và 14/15 loại theo pixel-AUROC, và thứ tự PatchCore, PaDiM,
Autoencoder giữ nguyên ở cả nhóm bề mặt lẫn nhóm vật thể (Bảng~\ref{tab:khao-sat}). Kết quả của PatchCore (0,983 và
0,980) sát số liệu của bài báo gốc (0,990 và 0,981)~\cite{rothTotalRecallIndustrial2022}; chênh lệch nhỏ ở
image-AUROC có thể do nhóm không cắt vùng giữa $224\times224$ như bài gốc. Pixel-AUROC của Autoencoder (0,843) cũng gần
mức 0,82 được trích trong bài báo PaDiM~\cite{defardPaDiMPatchDistribution2021}. Như vậy quy trình đánh giá của nhóm
đủ tin cậy để dùng cho các thực nghiệm sau.

\begin{table}[htbp]
  \centering
  \caption{Image-AUROC theo từng loại trong khảo sát sơ bộ}
  \label{tab:khao-sat-loai}
  \small
  \begin{tabular}{@{}lccc@{\hspace{1.2cm}}lccc@{}}
    \toprule
    Loại & AE & PaDiM & PatchCore & Loại & AE & PaDiM & PatchCore \\
    \midrule
    carpet  & 0,553 & 0,988 & 0,988 & hazelnut   & 0,923 & 0,720 & 1,000 \\
    grid    & 0,857 & 0,829 & 0,990 & metal\_nut & 0,574 & 0,957 & 0,997 \\
    leather & 0,807 & 0,999 & 1,000 & pill       & 0,603 & 0,824 & 0,942 \\
    tile    & 0,766 & 0,943 & 1,000 & screw      & 0,579 & 0,752 & 0,964 \\
    wood    & 0,938 & 0,966 & 0,987 & toothbrush & 0,958 & 0,772 & 0,917 \\
    bottle  & 0,926 & 0,998 & 1,000 & transistor & 0,807 & 0,943 & 0,995 \\
    cable   & 0,579 & 0,838 & 0,988 & zipper     & 0,776 & 0,927 & 0,977 \\
    capsule & 0,689 & 0,903 & 0,994 &            &       &       &       \\
    \bottomrule
  \end{tabular}
\end{table}

Giá trị trung bình che đi nhiều khác biệt giữa các loại (Bảng~\ref{tab:khao-sat-loai}). Ba loại được chọn cho thực
nghiệm chính, mỗi loại đại diện cho một kiểu điểm yếu. Screw là vật thể nhỏ xoay tự do, có image-AUROC thấp với cả
Autoencoder (0,579) và PaDiM (0,752); cùng với hazelnut, đây là nơi PaDiM kém PatchCore nhiều nhất, phù hợp với việc
thống kê theo vị trí bị ảnh hưởng khi vật thể không cố định. Grid là bề mặt lặp lại và không xoay, nhưng PaDiM vẫn
kém PatchCore 0,161; góc xoay không giải thích được khoảng cách này, nên cần kiểm tra ảnh hưởng của backbone.
Transistor là vật thể cố định, căn chỉnh tốt, có lỗi logic; cả PaDiM và PatchCore đều tốt, nên đây là loại đối
chứng.

Khảo sát có hai hạn chế mà các thực nghiệm sau đã khắc phục: PaDiM và PatchCore dùng backbone khác nhau nên không tách
được ảnh hưởng của cơ chế với ảnh hưởng của backbone, và mỗi phương pháp chỉ chạy một lần. Khi chuyển PaDiM sang
Wide-ResNet-50-2, image-AUROC trên grid tăng từ 0,829 lên 0,934, còn trên screw không tăng (0,752 so với 0,722). Khoảng
cách trên grid vì vậy chủ yếu do backbone, còn khoảng cách trên screw gắn với giả định căn chỉnh của PaDiM. Do cấu
hình khác, số liệu ở Bảng~\ref{tab:khao-sat-loai} chỉ dùng để chọn loại, không dùng để so sánh với các mục sau.
